% Generated draft: requires mathematical and source review.
\documentclass[11pt]{article}
\usepackage[margin=25mm]{geometry}
\usepackage{fontspec}
\setmainfont{Latin Modern Roman}
\usepackage{amsmath,amssymb,mathtools,unicode-math}
\setmathfont{Latin Modern Math}
\usepackage{enumitem,needspace,xurl,hyperref}
\setlength{\parindent}{0pt}
\setlength{\parskip}{5pt}
\setlength{\emergencystretch}{3em}
\newcommand{\N}{\mathbb{N}}
\newcommand{\Z}{\mathbb{Z}}
\newcommand{\Q}{\mathbb{Q}}
\newcommand{\R}{\mathbb{R}}
\newcommand{\C}{\mathbb{C}}
\newcommand{\F}{\mathbb{F}}
\newcommand{\A}{\mathbb{A}}
\newcommand{\PP}{\mathbb{P}}
\newcommand{\E}{\mathbb{E}}
\newcommand{\eps}{\varepsilon}
\newcommand{\dd}{\,\mathrm{d}}
\begin{document}
\section*{UnsolvedMath ID: 30005742}
\textbf{Source problem code:} catalogue entry
\subsection{Definitions and mathematical statement}
For terminology only, let \(k\) be a field and let \(A\) be a commutative unital Poisson \(k\)-algebra: its bracket \(\{-,-\}\) is a \(k\)-bilinear Lie bracket satisfying
\[
\{a,bc\}=\{a,b\}c+b\{a,c\}\qquad(a,b,c\in A).
\]
An ideal \(I\subseteq A\) is a Poisson ideal if \(\{A,I\}\subseteq I\). A prime Poisson ideal means a proper Poisson ideal that is prime as a commutative ideal. For an ideal \(J\subseteq A\), define its Poisson core by
\[
\operatorname{Pcore}(J)=\sum_{\substack{I\subseteq J\\ I\text{ a Poisson ideal}}}I.
\]
Under the usual convention, a Poisson primitive ideal is an ideal of the form \(\operatorname{Pcore}(\mathfrak m)\) for a maximal ideal \(\mathfrak m\) of \(A\). These definitions supply terminology; the excerpt does not specify its ground field, hypotheses on \(A\), or precise convention for Poisson primitivity.

For a Lie \(k\)-algebra \(\mathfrak g\), its symmetric algebra \(S_k(\mathfrak g)\) carries the Lie--Poisson bracket determined by \(\{x,y\}=[x,y]\) for \(x,y\in\mathfrak g\), extended by the Leibniz rule. The source writes \(S(\mathrm{Vir})\), with \(\mathrm{Vir}\) referring to the Virasoro algebra, but does not supply its ground field or presentation. An algebra \(A\) satisfies the ascending chain condition on prime Poisson ideals if every chain
\[
P_1\subseteq P_2\subseteq\cdots
\]
of prime Poisson ideals stabilizes: there is \(N\geq1\) such that \(P_n=P_N\) for all \(n\geq N\).

\textbf{Recoverable questions and unresolved scope.} The excerpt asks for a reasonable parametrization of an antecedently specified family of ideals, preferably yielding further information about each ideal, and for criteria distinguishing Poisson primitive ideals from other prime Poisson ideals. The antecedent of ``them'' and the ambient algebra or class of algebras are absent. Consequently, neither a quantified classification problem nor a required meaning of ``reasonable'' can be recovered precisely. The title suggests a connection with the Virasoro algebra, but does not determine the missing scope.

The excerpt further attributes to reference [4, Corollary 5.1] the assertion that \(S(\mathrm{Vir})\) satisfies the ascending chain condition on prime Poisson ideals, and asks ``What are the others?'' The class of objects quantified over by ``others'' is unspecified; this cannot safely be converted into a classification of algebras satisfying that chain condition or into a classification of ideals.

\textbf{Status qualification.} Although the supplied registry labels the record open, its statement-verification field says that the report presents these questions as motivation and subsequently states that the authors answer all of them. The same record marks the statement unrecoverable. Thus the supplied evidence does not establish a precise unresolved mathematical problem; the passages above preserve the recoverable questions without asserting that they remain open.
\subsection{Short English statement}
The excerpt asks for parametrizations and distinguishing criteria for Poisson primitive ideals, then mentions an ascending chain condition for \(S(\mathrm{Vir})\). Missing context leaves the intended families and the meaning of ``others'' unclear. The supplied open label conflicts with a verification note saying that the report's authors answer these motivational questions. This is a statement-only draft requiring human mathematical review.
\subsection{Sources}
\begin{itemize}
\item Supplied local UnsolvedMath registry record, permanent ID \(30005742\), title \emph{Poisson Ideals of the Virasoro Algebra}, source problem code \texttt{catalogue entry}, in the Oberwolfach Reports Open Problems collection. No proposer is supplied; the proposed year is recorded as 2023.
\item Supplied report citation: \emph{Mini-Workshop: Poisson and Poisson-type algebras} (2024), Oberwolfach Reports, pp.\ 2681--2715, DOI \texttt{10.4171/owr/2023/46}. The extracted passage includes page 2707. Supplied source URL: \url{https://doi.org/10.4171/owr/2023/46}. The record also labels this source \emph{Oberwolfach Report 46/2023}; the differing year designations are retained without resolution.
\item Supplied literature source: \emph{Oberwolfach Report 46/2023}, \url{https://doi.org/10.4171/owr/2023/46}. Literature checked at: \texttt{2026-08-22}. This is supplied dated curation, not an independently verified assessment. It labels the questions open and says the cited ascending-chain result for the Virasoro symmetric algebra does not settle the general problem. The supplied background explicitly requests independent verification of time-sensitive status and cited claims.
\item The excerpt cites [4, Corollary 5.1] for the ascending chain condition on prime Poisson ideals of \(S(\mathrm{Vir})\). The bibliographic entry for [4] is not supplied, so its authors, title, hypotheses, and precise conclusion cannot be identified from this record.
\item The supplied statement-verification note says the official report treats these questions as motivation and states that the authors answer all of them. This conflicts with the supplied open label and dated literature assessment. No independent source inspection is asserted here.
\end{itemize}
\end{document}
